\section{Scriptural Date and Location}
\label{sec:date-location}
{\small Each passage the Mass reads or sings from Scripture, in canonical order,
with what is dated, where it belongs, and the date the Scripture chronology
corpus gives for it. A passage whose parts are dated differently has a row for
each part; rows that share one answer name every element that uses it. The
Entrance Antiphon, both Communion antiphons and the Alleluia verse are placed at
the verses they adapt or excerpt; the three orations carry no biblical date.
The notes beneath give each passage's setting, attribution and dating
evidence.\par}
\medskip

\begin{chronologyorientationtable}
\chronologyorientationrow{row-29cbe0d234a47aa8}{Occasion and place unknown; the psalms' titles name none.}
\addlinespace
\chronologyorientationrow{row-6e9c8175463dbd35}{Place of composition unknown.}
\addlinespace
\chronologyorientationrow{row-d7a88409f6dff8dd}{Ascribed to David by its Latin and Greek title; David reigns at Hebron, then at Jerusalem. The reign dates David, not the psalm.}
\addlinespace
\chronologyorientationrow{row-b6be5603e767096a}{The house rebuilt after the captivity, named in the psalm's Latin and Greek title: the Temple at Jerusalem.}
\addlinespace
\chronologyorientationrow{row-7a3e8d424af7908c}{The Messiah's coming that the Fathers read in this psalm, taken as the Nativity at Bethlehem.}
\midrule
\chronologyorientationrow{row-2da3a3c2bfe61ef2}{Isaias, a citizen of Jerusalem, prophesies concerning Juda and Jerusalem.}
\addlinespace
\chronologyorientationrow{row-e535d564c1cee995}{An unnamed prophet, place not stated; his oracles bid the exiles go out of Babylon.}
\midrule
\chronologyorientationrow{row-a180ccaa931e671d}{Jesus answers the tribute question while teaching in the temple at Jerusalem.}
\addlinespace
\chronologyorientationrow{row-5396102e03d2c0da}{Matthew writes for Jewish readers; the place is not fixed (Eusebius: on leaving Palestine).}
\midrule
\chronologyorientationrow{row-4b690f3afb3478b9}{Mark writes at Rome for Roman Christians (Chrysostom names Egypt).}
\midrule
\chronologyorientationrow{row-11f6d60f0fa3c25e}{Paul writes from captivity at Rome (some say Caesarea) to the church at Philippi.}
\midrule
\chronologyorientationrow{row-d67f41bffc530c4a}{Paul writes from Corinth (an ancient subscription says Athens) to the church at Thessalonica.}
\end{chronologyorientationtable}

\subsection{Psalms 17 (16), 33 (32) and 96 (95): the antiphons and the psalm}
\label{sec:date-psalms}

None of the three psalms names an occasion in a way the chronology record can
date, except for Psalm 96 (95), whose Latin and Greek title does. The
``Historical setting'' row says so for the Entrance Antiphon's psalm and the
first Communion antiphon's: the Vulgate titles them ``The prayer of David''
and ``A psalm for David'' without naming an event or a place. The composition
date shown for all three is the bound that the introduction to the Psalms in
the \work{New American Bible, Revised Edition} sets for the whole Psalter,
within which, it holds, no individual psalm can be dated securely. It is not a
date for any of these three.

Psalm 33 (32) alone carries an attribution answer. Its title, \latin{Psalmus
David} in the Clementine and \textit{T\=o Daueid} in the Greek, ascribes it to
David, and the record dates David's reign in what the \work{Catholic
Encyclopedia}'s ``David'' (Corbett, vol.~4) calls ``the usual chronology'',
marked disputed because the article reports later datings beside it. The same
article places David's kingship first at Hebron and then at Jerusalem. The
reign is the reference point of the attribution and not a date of the psalm.
Psalm 17 (16) is also titled for David in the Vulgate, but the record holds no
attribution answer for it, and none is given.

Psalm 96 (95) is the responsorial psalm. Its title in the Clementine and in
Swete's Septuagint sets it ``when the house was built after the captivity'';
the Hebrew psalm and the \work{Nova Vulgata} carry no title, and the title is
not sung. The superscription row dates the setting the title names, the
raising of the Temple at Jerusalem after the return, on the figure the
\work{Catholic Encyclopedia}'s ``Temple of Jerusalem'' (Meistermann, vol.~14)
gives. It dates the title's setting, not the writing of the psalm. The
prophetic-referent row records the same encyclopedia's statement, in its
``Psalms'' (Drum, vol.~12), that the coming of the Messiah is foretold in this
group of psalms, and the chronology record's own step from that coming to the
Nativity; the two dates shown are those its articles ``General Chronology''
(Gerard) and ``Biblical Chronology'' (Howlett, both vol.~3) report, and the
record marks them disputed. They date the event the psalm was read as
foretelling, not the psalm.

\subsection{Isaiah 45:1, 4--6: the first reading}
\label{sec:date-isaiah}

The first reading has two answers, and they are not two dates for one thing.
The traditional attribution is to Isaias, the son of Amos (Is 1:1), whose
ministry the \work{Catholic Encyclopedia}'s ``Isaias'' (Souvay, vol.~8)
describes as lasting from the closing year of King Ozias to the latest
prophecies it admits; the article calls him a citizen, perhaps a native, of
Jerusalem, and the same paragraph reports that commentators differ on when he
was called. The row dates the prophet's ministry in that account and not the
oracle to Cyrus. The second answer is the critical placement of the oracles of
chapters 40--55: the introduction to Isaiah in the \work{New American Bible,
Revised Edition} attributes them to an anonymous prophet who spoke toward the
end of the Babylonian exile. That introduction gives the prophet no place; the
oracles themselves bid the exiles ``Come forth out of Babylon'' (48:20).
Neither answer is a date for the composition of 45:1--6, for which the record
holds none.

\subsection{Matthew 22:15--21: the Gospel}
\label{sec:date-matthew}

The ``Event'' row has no date: the chronology corpus holds none for the
question about tribute, and none is supplied here. Matthew sets it in the
temple at Jerusalem, where Jesus is teaching (21:23) after his entry into the
city (21:10), and before he leaves the temple (24:1); the \work{Catholic
Encyclopedia}'s ``Gospel of St.~Matthew'' (Jacquier, vol.~10) groups these
chapters under the heading ``In Jerusalem'', among them the questions
``concerning the tribute''. The composition dates shown are those the same
article reports from the ancient writers, who ``are at variance'', together
with the figures it gives where opinion is divided, and one from the
encyclopedia's ``The New Testament'' (Durand, vol.~14); all are marked
disputed. The article says that Matthew wrote for the Jews, and reports
Eusebius's account that he wrote in Hebrew when he was leaving Palestine; the
place of writing is not fixed. These are dates of the book, not of the saying.

\subsection{Mark 10:45: the second Communion antiphon}
\label{sec:date-mark}

The antiphon's source is dated only as part of Mark's Gospel. The figures
shown are those the \work{Catholic Encyclopedia}'s ``Gospel of Saint Mark''
(MacRory, vol.~9) reports from a manuscript subscription, the Paschal
Chronicle and Eusebius's Chronicle, together with the article's own conclusion,
a figure from its ``The New Testament'' (Durand), and the answer of the
Pontifical Biblical Commission of 1912 placing the Gospel after Matthew and
before the destruction of Jerusalem; all are marked disputed. MacRory holds that
the Gospel was written at Rome for Roman Christians, and reports, to reject it,
Chrysostom's statement that Mark wrote in Egypt.

\subsection{Philippians 2:15d, 16a: the Alleluia verse}
\label{sec:date-philippians}

The verse is dated as part of the letter. The first figure is the year the
\work{Catholic Encyclopedia}'s ``St.~Paul'' (Prat, vol.~11) gives the letters
of the captivity; the second is the range its ``Epistle to the Philippians''
(Vander Heeren, vol.~12) gives for Paul's captivity at Rome. Both are marked
disputed. Vander Heeren reports that some have held the letter to have been
written at Caesarea, and holds that it was written at Rome; the letter itself
names Paul's bonds (1:13) and greets the Philippians from ``Caesar's household''
(4:22). The verse narrates no event.

\subsection{1 Thessalonians 1:1--5b: the second reading}
\label{sec:date-thessalonians}

The figures shown are those of the \work{Catholic Encyclopedia}'s ``St.~Paul''
(Prat) and ``Epistles to the Thessalonians'' (Drum, vol.~14), the second of
which gives its years ``according to the commonly received scheme of Pauline
chronology''; both are marked disputed. Drum reports that the received Greek
text and several uncials subscribe the letters ``written from Athens'', and
holds from the letter itself (3:6) and from Acts (18:5) that they were written
during Paul's first stay in Corinth. The letter is addressed ``to the church of
the Thessalonians'' (1:1), whose founding Acts narrates (17:1--10).
